\clearpage\pagestyle{otMain}\setlength{\otBlocEspacement}{2.5mm}\fontsize{11}{15}\selectfont\corrigesChapitre{13}{Les angles}
\begin{blocLivre}{cor-13-1}\begin{corrige}{1}{Nommer et décrire}\demarche{Repérer le sommet d’abord.}
C’est l’angle $\widehat{ABC}$ : sommet $\res B$, côtés passant par A et par C. La lettre du sommet s’écrit au milieu.
\label{fig-cor-13-1}\figureGeo[corrige]{ang-1}
\end{corrige}\end{blocLivre}
\begin{blocLivre}{cor-13-2}\begin{corrige}{2}{Comparer au calque}\demarche{Superposer sommet et un côté.}
Le calque superposé donne $\res{\widehat B<\widehat A<\widehat C}$. Les longueurs des côtés dessinés ne comptent pas.
\label{fig-cor-13-2}\figureGeo[corrige]{ang-2}
\end{corrige}\end{blocLivre}
\begin{blocLivre}{cor-13-3}\begin{corrige}{3}{Reproduire un angle}\demarche{Garder le calque bien calé pendant le tracé.}
On décalque l’angle, on pose le calque sur la demi-droite, sommet sur le point marqué, puis on trace le second côté. La reproduction figure en vert.
\label{fig-cor-13-3}\figureGeo[corrige]{ang-3}
\end{corrige}\end{blocLivre}
\begin{blocLivre}{cor-13-4}\begin{corrige}{4}{La chasse aux angles droits}\demarche{Tester chaque sommet à l’équerre.}
L’équerre s’ajuste en $\res A$ et en $\res D$ : deux angles droits, codés d’un petit carré. Les deux autres angles ne conviennent pas.
\label{fig-cor-13-4}\figureGeo[corrige]{ang-4}
\end{corrige}\end{blocLivre}
\begin{blocLivre}{cor-13-5}\begin{corrige}{5}{90 degrés}\demarche{Composer à partir de 90.}
L’angle plat vaut deux droits : $90+90=\res{180^\circ}$. La moitié d’un droit vaut $90:2=\res{45^\circ}$.
\end{corrige}\end{blocLivre}
\begin{blocLivre}{cor-13-6}\begin{corrige}{6}{La somme de deux angles}\demarche{Juxtaposer sans chevaucher.}
On reporte le second angle contre un côté du premier, sommets superposés. La somme, en vert, est $\res{\text{plus ouverte qu’un angle droit}}$ : l’équerre ne la couvre pas ; ici $55+45=100>90$.
\label{fig-cor-13-6}\figureGeo[corrige]{ang-6}
\end{corrige}\end{blocLivre}
\begin{blocLivre}{cor-13-7}\begin{corrige}{7}{Le double}\demarche{Reporter deux fois avec le calque.}
On reporte l’angle deux fois côte à côte : le double figure en vert et mesure $35\times2=\res{70^\circ}$, encore plus petit qu’un angle droit.
\label{fig-cor-13-7}\figureGeo[corrige]{ang-7}
\end{corrige}\end{blocLivre}
\begin{blocLivre}{cor-13-8}\begin{corrige}{8}{La moitié par pliage}\demarche{Amener côté sur côté par pliage.}
On plie le calque de l’angle en amenant un côté exactement sur l’autre : le pli donne la moitié, en vert. Pour un angle droit : $\res{45^\circ}$.
\label{fig-cor-13-8}\figureGeo[corrige]{ang-8}
\end{corrige}\end{blocLivre}
\begin{blocLivre}{cor-13-9}\begin{corrige}{9}{Des mesures sans rapporteur}\demarche{Composer avec 90 et 45.}
$90+45=\res{135^\circ}$ ; $45\times3=\res{135^\circ}$ aussi ; $180-90=\res{90^\circ}$ : un angle droit.
\end{corrige}\end{blocLivre}
\begin{blocLivre}{cor-13-10}\begin{corrige}{10}{Ranger des angles}\demarche{Situer chaque angle par rapport à 90°.}
$\widehat D$ est plus petit qu’un droit, $\widehat A$ est droit, $\widehat C$ dépasse le droit, $\widehat B$ est presque plat : $\res{\widehat D<\widehat A<\widehat C<\widehat B}$.
\label{fig-cor-13-10}\figureGeo[corrige]{ang-10}
\end{corrige}\end{blocLivre}
\begin{blocLivre}{cor-13-11}\begin{corrige}{11}{Les aiguilles de l’horloge}\demarche{Imaginer le cadran et les quarts de tour.}
À $\res{3\ \text h}$ et à $\res{9\ \text h}$, les aiguilles forment un angle droit. À 6 h, elles sont opposées : un angle $\res{\text{plat}}$, de 180 degrés.
\label{fig-cor-13-11}\figureGeo[corrige]{ang-11}
\end{corrige}\end{blocLivre}
\begin{blocLivre}{cor-13-12}\begin{corrige}{12}{L’éventail}\demarche{Partager 180 en parts égales.}
$180:4=\res{45^\circ}$ : chaque secteur vaut la moitié d’un angle droit. Contrôle : $4\times45=180$.
\label{fig-cor-13-12}\figureGeo[corrige]{ang-13}
\end{corrige}\end{blocLivre}
\begin{blocLivre}{cor-13-13}\begin{corrige}{13}{Le toit de la cabane}\demarche{Partager l’angle droit en deux.}
L’angle droit se partage en deux angles égaux : $90:2=\res{45^\circ}$ chacun. Le pliage du calque confirmerait l’égalité.
\label{fig-cor-13-13}\figureGeo[corrige]{ang-14}
\end{corrige}\end{blocLivre}
\begin{blocLivre}{cor-13-14}\begin{corrige}{14}{Des côtés trompeurs}\demarche{Comparer les écartements au calque.}
L’ouverture seule compte : le calque montre $\res{\widehat E<\widehat F}$ malgré ses longs côtés dessinés. Un angle n’est pas une surface ni une longueur.
\label{fig-cor-13-14}\figureGeo[corrige]{ang-15}
\end{corrige}\end{blocLivre}
\begin{blocLivre}{cor-13-15}\begin{corrige}{15}{Le programme d’angles}\demarche{Additionner les mesures reportées.}
L’angle total juxtapose $90^\circ$ et $45^\circ$ : il mesure $90+45=\res{135^\circ}$. La construction, au calque et à l’équerre, figure en vert.
\label{fig-cor-13-15}\figureGeo[corrige]{ang-9}
\end{corrige}\end{blocLivre}
